% The spectral census of the icosahedral design at (6,3): all twenty triples,
% the sign grid that decides them, and the invariants they share.  There is no
% projection -- the plate is three panels of tabular material -- so the
% coordinate derivation below is the layout metric.  x = y = 1cm throughout.
%
% LAYOUT.
%   left, the sign grid   x 0.55 .. 4.15   cell 0.60, six rows and six columns,
%                         cell (i,j) spanning [0.55+0.60j, 1.15+0.60j] x
%                         [-0.45-0.60i, -1.05-0.60i], centre
%                         (0.85+0.60j, -0.75-0.60i); column heads at
%                         (0.85+0.60j, -0.20), row heads at (0.30, -0.75-0.60i)
%   middle, the census    x 4.55 .. 7.25   two columns centred at 5.10 and 6.70,
%                         heads at -0.20, ten rows at -0.68 - 0.34r, r = 0..9
%   right, the invariants x 7.55 .. 15.51  labels at 7.60, values centred at
%                         11.15 and 14.15; head row at -0.20, rule at -0.42,
%                         nine rows at -0.68 - 0.40r, r = 0..8, and a second
%                         rule at -3.66 setting the verdict row off from the
%                         eight invariants above it
%   footer                four lines at -4.40, -4.72, -5.04, -5.36, from x = 0
% Panel titles sit at y = 0.15.  The widest cell in the two value columns is a
% spectrum, at 2.72cm, and the widest row label is 1.93cm, so the three columns
% of the right panel need 1.95 + 2.72 + 2.72 with two gaps; the plate is 15.51cm
% wide, inside the 16.20cm text block.
%
% THE DESIGN, exactly as in the appendix.  rho = 3/sqrt5, sigma^2 = (3-rho)/2,
% mu^2 = (3+rho)/2, so sigma^2 + mu^2 = 3 and sigma^2 mu^2 = (9 - rho^2)/4
% = (9 - 9/5)/4 = 9/5 = rho^2, whence sigma mu = rho.  The six atoms are
%   g_0 = ( sigma, mu, 0)    g_2 = (0,  sigma, mu)    g_4 = (mu, 0,  sigma)
%   g_1 = (-sigma, mu, 0)    g_3 = (0, -sigma, mu)    g_5 = (mu, 0, -sigma)
% at weight 1/6.  Each coordinate receives 2(sigma^2 + mu^2) = 6 from the six
% outer products and each off-diagonal cancels in pairs, so
% sum_c t_c g_c g_c^T = I_3; every leverage is 3, and every distinct pairing is
% +-sigma mu or +-(mu^2 - sigma^2), both of which are +-rho, so
% <g_c, g_d>^2 = rho^2 = 9/5.
%
% THE SIGN GRID, computed in exact algebraic arithmetic from those six vectors.
% Entry (c,d) is <g_c, g_d>/rho, which is +1 or -1 off the diagonal:
%        d = 0   1   2   3   4   5
%   c = 0    3   +   +   -   +   +
%   c = 1    +   3   +   -   -   -
%   c = 2    +   +   3   +   +   -
%   c = 3    -   -   +   3   +   -
%   c = 4    +   -   +   +   3   +
%   c = 5    +   -   -   -   +   3
% Six of the fifteen pairs are negative: 03, 13, 14, 15, 25, 35.  Replacing
% g_c by -g_c flips the five pairings at c; a triple containing c has two of
% its three pairings flipped, so no product tau changes sign, and a triple
% missing c is untouched.  The grid is therefore fixed only up to that
% switching, and every verdict below is not.
%
% THE CENSUS.  Writing tau for the product of the three pairings of a triple,
% tau = +-rho^3 always, and the twenty triples split ten and ten:
%   tau = +rho^3   012 013 024 035 045 125 134 145 234 235
%   tau = -rho^3   014 015 023 025 034 123 124 135 245 345
% Both classes have Gram diagonal 3 and every off-diagonal square 9/5, so
%   tr Gram_C = 9   and   sum_{c<d in C} <g_c,g_d>^2 = 3 (9/5) = 27/5
% at all twenty, while
%   det Gram_C = 27 + 2 tau - 3(27/5) = 54/5 + 2 tau,
%   gamma(C) = 4 det Gram_C - (tr Gram_C - 1)^2 = 8 tau - 104/5.
% A triple with all three pairings equal to +rho has S_C = (3-rho) I_3 + rho J_3,
% of spectrum {3-rho, 3-rho, 3+2rho}; conjugating by a diagonal sign matrix
% changes the pattern but not the spectrum, and every pattern with tau > 0 is
% reached that way, so that spectrum holds at all ten.  The other ten carry
% {3-2rho, 3+rho, 3+rho}.  Both have trace 9, as they must.  Numerically the
% two least eigenvalues are 1.658359... and 0.316718..., and the two values of
% gamma are -1.480372... and -40.119627..., all four negative in gamma, but the
% plate prints no decimal: 3 - rho > 1 because rho^2 = 9/5 < 4, and 3 - 2rho < 1
% because rho^2 > 1.
%
% THE BRIDGE TO THE MECHANIZED CRITERION.  The Lean criterion is stated in the
% pairings normalized by the heavy excesses ell_c - 1 = 2, so the normalized
% pairing is <g_c,g_d>/2 and its square is (9/5)/4 = 9/20.  Its product over a
% triple has modulus (9/20)^{3/2} = 27/(40 sqrt5) = 27 sqrt5 / 200, and
% 27 sqrt5 = 60.373... exceeds 35, so that modulus exceeds 7/40 = 35/200.  The
% criterion "7/40 <= the product" is therefore the statement that the product is
% positive, which is the sign of tau.
%
% DISCLOSED.  The split ten and ten is the one number here with no theorem
% behind it; the development records it in a docstring, not a statement, and
% Appendix B says so.  The design, the leverage 3, the single squared pairing
% 9/5, the criterion in the normalized pairings, the strict domination of
% {0,2,4} and the two-sided margin are all kernel-checked; the sign grid, the
% two ten-element lists, the two spectra and the arithmetic 27 sqrt5 > 35 are
% computed here.
\begin{tikzpicture}[
  x=1cm, y=1cm,
  cellpos/.style={fill=black!10, draw=black!45, line width=0.25pt},
  cellneg/.style={fill=black!45, draw=black!45, line width=0.25pt},
  celldiag/.style={fill=black!25, draw=black!45, line width=0.25pt},
  sgn/.style={font=\scriptsize, inner sep=0pt},
  head/.style={font=\scriptsize, inner sep=0pt},
  title/.style={font=\footnotesize, anchor=west, inner sep=0pt},
  cell/.style={font=\scriptsize, inner sep=0pt},
  lab/.style={font=\scriptsize, anchor=west, inner sep=0pt},
  note/.style={font=\scriptsize, anchor=west, inner sep=0pt},
  hair/.style={line width=0.25pt, black!45}]

% ========================= left: the sign grid ==========================
\node[title] at (0.00,0.15) {the pairing signs};
\foreach \jj in {0,...,5}{%
  \node[head] at ({0.85+0.60*\jj},-0.20) {\jj};
  \node[head] at (0.30,{-0.75-0.60*\jj}) {\jj};}
% row c = 0 : 3 + + - + +
\path[celldiag] (0.55,-0.45) rectangle (1.15,-1.05);
\foreach \jj in {1,2,4,5}{\path[cellpos]
  ({0.55+0.60*\jj},-0.45) rectangle ({1.15+0.60*\jj},-1.05);}
\path[cellneg] (2.35,-0.45) rectangle (2.95,-1.05);
% row c = 1 : + 3 + - - -
\path[celldiag] (1.15,-1.05) rectangle (1.75,-1.65);
\foreach \jj in {0,2}{\path[cellpos]
  ({0.55+0.60*\jj},-1.05) rectangle ({1.15+0.60*\jj},-1.65);}
\foreach \jj in {3,4,5}{\path[cellneg]
  ({0.55+0.60*\jj},-1.05) rectangle ({1.15+0.60*\jj},-1.65);}
% row c = 2 : + + 3 + + -
\path[celldiag] (1.75,-1.65) rectangle (2.35,-2.25);
\foreach \jj in {0,1,3,4}{\path[cellpos]
  ({0.55+0.60*\jj},-1.65) rectangle ({1.15+0.60*\jj},-2.25);}
\path[cellneg] (3.55,-1.65) rectangle (4.15,-2.25);
% row c = 3 : - - + 3 + -
\path[celldiag] (2.35,-2.25) rectangle (2.95,-2.85);
\foreach \jj in {0,1,5}{\path[cellneg]
  ({0.55+0.60*\jj},-2.25) rectangle ({1.15+0.60*\jj},-2.85);}
\foreach \jj in {2,4}{\path[cellpos]
  ({0.55+0.60*\jj},-2.25) rectangle ({1.15+0.60*\jj},-2.85);}
% row c = 4 : + - + + 3 +
\path[celldiag] (2.95,-2.85) rectangle (3.55,-3.45);
\foreach \jj in {0,2,3,5}{\path[cellpos]
  ({0.55+0.60*\jj},-2.85) rectangle ({1.15+0.60*\jj},-3.45);}
\path[cellneg] (1.15,-2.85) rectangle (1.75,-3.45);
% row c = 5 : + - - - + 3
\path[celldiag] (3.55,-3.45) rectangle (4.15,-4.05);
\foreach \jj in {0,4}{\path[cellpos]
  ({0.55+0.60*\jj},-3.45) rectangle ({1.15+0.60*\jj},-4.05);}
\foreach \jj in {1,2,3}{\path[cellneg]
  ({0.55+0.60*\jj},-3.45) rectangle ({1.15+0.60*\jj},-4.05);}
% the glyphs
\foreach \ii/\ro in {0/{3,+,+,-,+,+}, 1/{+,3,+,-,-,-}, 2/{+,+,3,+,+,-},
                     3/{-,-,+,3,+,-}, 4/{+,-,+,+,3,+}, 5/{+,-,-,-,+,3}}{%
  \foreach \gl [count=\jj from 0] in \ro {%
    \node[sgn] at ({0.85+0.60*\jj},{-0.75-0.60*\ii}) {$\gl$};}}

% ========================== middle: the census ==========================
\node[title] at (4.55,0.15) {the twenty triples};
\node[cell] at (5.10,-0.20) {$\tau=+\rho^{3}$};
\node[cell] at (6.70,-0.20) {$\tau=-\rho^{3}$};
\draw[hair] (4.55,-0.42) -- (7.25,-0.42);
\foreach \tl [count=\rr from 0] in {012,013,024,035,045,125,134,145,234,235}
  {\node[cell] at (5.10,{-0.68-0.34*\rr}) {$\tl$};}
\foreach \tl [count=\rr from 0] in {014,015,023,025,034,123,124,135,245,345}
  {\node[cell] at (6.70,{-0.68-0.34*\rr}) {$\tl$};}

% ======================= right: the invariants ==========================
\node[title] at (7.55,0.15) {the invariants};
\node[cell] at (11.15,-0.20) {$\tau=+\rho^{3}$};
\node[cell] at (14.15,-0.20) {$\tau=-\rho^{3}$};
\draw[hair] (7.55,-0.42) -- (15.51,-0.42);
\foreach \rl/\va/\vb [count=\rr from 0] in {%
    {triples}/{$10$}/{$10$},
    {$\tr\Gram_C$}/{$9$}/{$9$},
    {$\sum_{c<d}\langle g_c,g_d\rangle^{2}$}/{$27/5$}/{$27/5$},
    {$\lvert\tau\rvert$}/{$\rho^{3}$}/{$\rho^{3}$},
    {$\det\Gram_C$}/{$54/5+2\rho^{3}$}/{$54/5-2\rho^{3}$},
    {$\operatorname{spec}\SC$}/{$\{3-\rho,3-\rho,3+2\rho\}$}/%
      {$\{3-2\rho,3+\rho,3+\rho\}$},
    {$\lmin(\SC)$}/{$3-\rho$}/{$3-2\rho$},
    {$\gamma(C)$}/{$8\rho^{3}-104/5$}/{$-8\rho^{3}-104/5$},
    {verdict}/{dominates}/{fails}}{%
  \node[lab]  at (7.60,{-0.68-0.40*\rr}) {\rl};
  \node[cell] at (11.15,{-0.68-0.40*\rr}) {\va};
  \node[cell] at (14.15,{-0.68-0.40*\rr}) {\vb};}
\draw[hair] (7.55,-3.66) -- (15.51,-3.66);

% ============================== footer =================================
\node[note] at (0,-4.40)
  {$\rho=3/\sqrt5$ and $\rho^{2}=9/5$.  The cell $(c,d)$ carries the sign of
   $\langle g_c,g_d\rangle/\rho$; the diagonal carries the leverage $3$.};
\node[note] at (0,-4.72)
  {Replacing $g_c$ by $-g_c$ flips the five pairings at $c$ and no product
   $\tau$, so the grid is fixed only up to that switching.};
\node[note] at (0,-5.04)
  {Normalized by the excesses $\ell_c-1=2$ the pairing has square $9/20$, so
   the product has modulus $27\sqrt5/200>7/40$:};
\node[note] at (0,-5.36)
  {domination is the sign of $\tau$.  And $3-\rho>1$ since $\rho^{2}<4$, while
   $3-2\rho<1$ since $\rho^{2}>1$.};

\end{tikzpicture}
